
\chapter{Introduction and Motivation}\label{chap:intro}

Historically, the development of efficient continuous optimization algorithms has relied on mathematical derivations and extensive manual trial-and-error by human experts over several decades \cite{Nelder1965, Hansen2001}. This effort has produced a wide range of general-purpose solvers, from classical gradient-based methods such as Gradient Descent and Quasi-Newton algorithms to widely used derivative-free heuristics, most notably the Nelder-Mead simplex algorithm \cite{Nelder1965} and the Covariance Matrix Adaptation Evolution Strategy (CMA-ES) \cite{Hansen2001}. However, the diversity of real-world fitness landscapes means that no single general-purpose solver can achieve optimal performance across all problem instances without some form of specialized adaptation. Customizing algorithmic logic or tuning parameters for each application therefore requires substantial engineering effort and does not scale well across distinct optimization problems \cite{Burke2013}.

This challenge becomes more pronounced when optimization algorithms are deployed in non-deterministic, noisy environments. Traditional ranking-based heuristics such as CMA-ES can perform poorly in noisy settings, where statistical fluctuations may lead the search toward false improvements and require additional statistical mechanisms to maintain reliable convergence \cite{Hansen2009Noisy}. Designing a robust mathematical algorithm to handle this uncertainty is difficult because a fundamentally poor candidate solution may temporarily obtain a favorable fitness value due to a statistical anomaly, causing the search process to adjust its parameters toward an incorrect direction \cite{Jin2005}. The problem is particularly pronounced for classical gradient-based algorithms. Since these methods rely on analytical or finite-difference derivatives, even relatively small amounts of statistical noise can corrupt gradient estimates and substantially affect the resulting search direction \cite{Sarafian2024}.

To automate the discovery of specialized search strategies and reduce the need for manual algorithm design, techniques such as hyperparameter tuning, automated algorithm selection, and traditional Automated Algorithm Design (AAD) were developed long before the advent of deep learning. More recently, AAD has gained renewed interest through the use of Large Language Models (LLMs). Rather than relying entirely on rigid, predefined search templates, modern AAD frameworks incorporate foundation models into structured, closed-loop environments that can write, test, and refine executable source code \cite{Romera2023}.

To implement such a code-generation loop, this research leverages the Large Language Model Evolutionary Algorithm (LLaMEA) framework \cite{LLaMEA2024}. The pipeline prompts the generative model to produce fully functional Python optimization classes, executes the generated code within a secure validation sandbox against standard mathematical benchmarks, and feeds execution metrics and compiler tracebacks back into the LLM's context to guide subsequent mutation and debugging. Existing work has shown that frameworks such as LLaMEA can produce competitive optimization results under deterministic evaluation conditions \cite{LLaMEA2024}. However, their ability to adapt generated code structures to non-deterministic conditions remains largely unexplored. This thesis addresses this gap by evaluating the performance and adaptive behavior of LLaMEA across continuous black-box optimization landscapes affected by persistent statistical noise, investigating whether closed-loop evolutionary feedback can lead language models to synthesize noise-resilient search strategies.

This thesis establishes a benchmark-driven evaluation of Automated Algorithm Design under uncertainty. A controlled stochastic layer is constructed over the continuous Black-Box Optimization Benchmark (BBOB) suite \cite{Hansen2010}, which is widely used for evaluating single-objective continuous optimization algorithms. By introducing parameter-controlled statistical noise into the benchmark functions, the evaluation environment provides a repeatable setting in which to examine the adaptive behavior of the LLaMEA framework. This configuration allows the study to focus on whether evolutionary pressure from code evaluation under stochastic conditions can lead an LLM to discover and generate specialized Python optimization heuristics for handling uncertainty.

The core contribution of this thesis is organized into three methodological phases. First, we formulate a non-deterministic evaluation sandbox by constructing a dynamic wrapper layer that applies controlled statistical noise profiles to standard benchmark problems. Second, we evaluate the ability of the generative model to modify Python-based metaheuristic class structures when exposed to these corrupted evaluations over multiple evolutionary cycles. Finally, we conduct a cross-environment robustness evaluation in which noise-trained champion models are compared against equivalent models evolved under clean, deterministic conditions, allowing the resulting algorithmic differences to be examined under the same evaluation settings.

The remainder of this thesis proposal is organized as follows. Chapter \ref{chap:research} reviews current work in black-box optimization, stochastic noise mitigation, and generative code-based algorithm design. Chapter \ref{chap:objectives} presents the research questions and formalizes the three experimental objectives, while the subsequent chapters describe the software pipeline architecture, scope and limitations, and provisional timeline for the final implementation.
